\documentclass[10pt,a4paper]{amsart}
\usepackage[margin=19mm]{geometry}
\usepackage{amsmath,amssymb,mathtools}
\usepackage[scr=rsfs,cal=euler]{mathalfa}
\usepackage{xfrac}
\usepackage[unicode,hidelinks,bookmarks=false]{hyperref}
\newcommand{\ie}{i.e.\ }
\newcommand{\cf}{cf.\ }
\newcommand{\resp}{respectively\ }
\newcommand{\Subsection}{\subsection}
\providecommand{\nobreakdash}{\nobreak\hbox{-}}
\setlength{\emergencystretch}{2em}
\multlinegap0pt
\usepackage{bm}
\usepackage{bbm}
\usepackage{dsfont}
\usepackage{xspace}
\usepackage{cancel}
\usepackage{mathtools}
\usepackage{amsthm}
\makeatletter
\@ifundefined{theorem}{\newtheorem{theorem}{Theorem}}{}
\@ifundefined{corollary}{\newtheorem{corollary}[theorem]{Corollary}}{}
\makeatother
\newcommand{\B}{\mathbb B}
\newcommand{\C}{\mathbb C}
\newcommand{\D}{\mathbb D}
\newcommand{\R}{\mathbb R}
\newcommand{\Sph}{\mathbb S}
\newcommand{\T}{\mathbb T}
\newcommand{\norm}[1]{\left\|#1\right\|}
\newcommand{\re}{\operatorname{Re}}
\title[The Schwarz lemma for holomorphic and minimal disks at the b]{The Schwarz lemma for holomorphic and minimal disks at the boundary}
\author{David Kalaj}
\date{Source: arXiv:2509.09471v4 (CC BY 4.0)}
\begin{document}
\maketitle

\noindent\emph{Source and licence.} The Schwarz lemma for holomorphic and minimal disks at the boundary. Version arXiv:2509.09471v4 (CC BY 4.0), distributed under the Creative Commons Attribution 4.0 International licence, as recorded in the arXiv licence metadata for this exact version and shown on the abstract page https://arxiv.org/abs/2509.09471v4. Full rights notice: licenses/arxiv-2509.09471-CC-BY-4.0.txt.\par\smallskip

\noindent\textit{Editorial scope.} Counted unit: the Corollary whose source label is \texttt{cor:hemisphere-lipschitz} in the article (source0.tex lines 920--929, source label \texttt{cor:hemisphere-lipschitz}), delivered together with the article's own complete original proof of it (source0.tex lines 931--965, one \texttt{proof} environment of 35 lines). No mathematics is written, repaired or re-derived: statement and proof are the article's own text line for line. Required context retained uncounted: th:minimal-main (lines 151--171); eq:weierstrass-data (lines 901--908); eq:radial-factor (lines 915--917). Every definition, display and label of those lines is retained unchanged. Omitted from this package and not redistributed: 1--150, 172--900, 909--914, 918--919, 930--930, 966--1051. Nothing inside a retained excerpt is omitted or altered: every retained body is delivered exactly as the article's lines are written, so no omission receipt of any kind accompanies this package. Citations: the retained excerpts carry no citation command at all, so no bibliography entry of the article is transcribed and no reference is redistributed. Reference adaptation: none was needed and none was made; every reference inside the delivered unit resolves inside it, so no unresolved reference or citation is produced. Printed slips kept literal: no printed word, symbol, hypothesis or number of the counted statement or of its proof was altered. Layout and class: the article's own class is \texttt{amsart} with packages including geometry, fontenc, inputenc, lmodern, amsmath, amssymb, amsthm, mathtools, mathrsfs, hyperref; the wrapper is the lane's pinned \texttt{amsart} editorial wrapper at 10pt on A4, which restates the theorem-like environments the retained text needs and adds the article's own macro definitions that the retained text needs (\texttt{\textbackslash B}, \texttt{\textbackslash C}, \texttt{\textbackslash D}, \texttt{\textbackslash R}, \texttt{\textbackslash Sph}, \texttt{\textbackslash T}, \texttt{\textbackslash norm}, \texttt{\textbackslash re}), with extra wrapper packages (bm, bbm, dsfont, xspace, cancel, mathtools). The delivered numbering is the unit's own. Assets: no retained span contains a figure environment, a graphics-inclusion command, a PDF-page inclusion, a table or a TikZ picture; no image file is copied into this package, the package declares no asset dependency and no image file is redistributed with it. Licence: version arXiv:2509.09471v4 of this article is distributed under the Creative Commons Attribution 4.0 International licence, as recorded in the arXiv licence metadata for this exact version and shown on the abstract page https://arxiv.org/abs/2509.09471v4. A full rights notice is delivered at licenses/arxiv-2509.09471-CC-BY-4.0.txt. Characters: every retained excerpt is pure ASCII, so the delivered engine, XeLaTeX with its default Latin Modern text font, reports no missing character.
\begin{theorem}[Boundary Schwarz lemma for conformal minimal disks]\label{th:minimal-main}
Let \(n\ge3\), and let \(F:\D\to\B^n\subset\R^n\) be a conformal minimal immersion. Assume that \(F\) has a boundary value \(F(\zeta)\in\Sph^{n-1}\) at some \(\zeta\in\T\), and that \(F\) extends differentiably to \(\zeta\) so that \(dF(\zeta)\) is defined. Then
\begin{equation}\label{eq:minimal-main-bound}
        \norm{dF(\zeta)}
        \ge
        \frac{1-\norm{F(0)}}{1+\norm{F(0)}}.
\end{equation}
If equality holds, then \(F(\D)\) is a totally geodesic linear disk. More precisely, there is a two-dimensional linear subspace \(L\subset\R^n\) such that
\[
        F(\D)=L\cap\B^n.
\]
In addition, if \(F(0)\ne0\), then
\[
        F(\zeta)=\frac{F(0)}{\norm{F(0)}}.
\]
Conversely, after identifying such an \(L\) with \(\C\), equality is attained, up to rotations of the source and of \(L\), by the disk automorphisms
\[
        z\longmapsto \frac{z+a}{1+az},
        \qquad 0\le a<1.
\]
\end{theorem}

\begin{equation}\label{eq:weierstrass-data}
        F(z)=\re\int^z
        \left(
        \frac12(1-q^2)p,
        \frac{i}{2}(1+q^2)p,
        qp
        \right).
\end{equation}

\begin{equation}\label{eq:radial-factor}
        |F_r(z)|=\norm{dF(z)}=\frac12 |p(z)|(1+|q(z)|^2),
\end{equation}

\begin{corollary}\label{cor:hemisphere-lipschitz}
Let \(F:\D\to\B^3\) be a conformal minimal embedding whose unit normals belong to an open hemisphere. Assume that \(F\) extends differentiably to \(\T\), that \(F(\T)\subset\Sph^2\), and that the Weierstrass data in \eqref{eq:weierstrass-data} extend continuously to \(\overline\D\). Then
\begin{equation}\label{eq:interior-lower-bound}
        \norm{dF(z)}
        \ge
        \frac12\frac{1-\norm{F(0)}}{1+\norm{F(0)}},
        \qquad z\in\D.
\end{equation}
Consequently, \(F^{-1}:F(\D)\to\D\) is Lipschitz continuous with respect to the intrinsic metric on \(F(\D)\) and the Euclidean metric on \(\D\).
\end{corollary}

\begin{proof}
Theorem~\ref{th:minimal-main} gives, for every \(\zeta\in\T\),
\[
        \norm{dF(\zeta)}
        \ge
        \frac{1-\norm{F(0)}}{1+\norm{F(0)}}.
\]
Since \(F\) is conformal, this norm is the conformal factor. In particular, by \eqref{eq:radial-factor},
\[
        \norm{dF(z)}=|F_r(z)|=\frac12 |p(z)|(1+|q(z)|^2).
\]
Since \(|q|<1\), we have \(\frac12(1+|q|^2)\le1\). The boundary estimate therefore gives
\[
        |p(\zeta)|
        \ge
        \frac{1-\norm{F(0)}}{1+\norm{F(0)}}.
\]
The immersion condition gives \(p\ne0\) in \(\D\). Applying the maximum principle to \(1/p\) yields the same lower bound for \(|p|\) throughout \(\D\). Hence \eqref{eq:interior-lower-bound} follows from \eqref{eq:radial-factor}, because \(\frac12(1+|q|^2)\ge\frac12\).

If \(\gamma\) is a \(C^1\) curve in \(\D\), then
\[
        \operatorname{length}_{F(\D)}(F\circ\gamma)
        \ge
        \frac12\frac{1-\norm{F(0)}}{1+\norm{F(0)}}
        \operatorname{length}_{\D}(\gamma).
\]
Taking infima over curves joining two points gives
\[
        |F^{-1}(x)-F^{-1}(y)|
        \le
        2\frac{1+\norm{F(0)}}{1-\norm{F(0)}}
        d_{F(\D)}(x,y),
\]
where \(d_{F(\D)}\) is the intrinsic distance on the minimal disk. This proves the Lipschitz assertion.
\end{proof}

\end{document}
